\documentclass[FM,DP]{tulthesis}

\usepackage{polyglossia}
\setdefaultlanguage{czech}
\usepackage{fontspec}
\usepackage{xunicode}
\usepackage{xltxtra}
\setsansfont[Mapping=tex-text,BoldFont={* Bold},Numbers=OldStyle]{Myriad Pro}
\usepackage{hyperref}
\hypersetup{colorlinks=true, linkcolor=tul, urlcolor=tul, citecolor=tul}
\usepackage{graphicx}
\usepackage{listings}
\usepackage[toc,page]{appendix}
\usepackage{amsmath}
\usepackage{amssymb}
\newcommand{\argument}[1]{{\ttfamily\color{\tulcolor}#1}}
\newcommand{\prikaz}[1]{\argument{\textbackslash #1}}
\newenvironment{myquote}{\begin{list}{}{\setlength\leftmargin\parindent}\item[]}{\end{list}}
\newenvironment{listing}{\begin{myquote}\color{\tulcolor}}{\end{myquote}}
\sloppy


%%%%%%%%%%%%%%%%%%%%%%%%%%%%
\TULtitle{Vyhledávání jako služba}{Search as a Service}
\TULprogramme{Aplikovaná informatika}{Applied Informatics}
\TULbranch{Informační systémy a technologie}{Information Technologies}
\TULauthor{Bc. Luděk Veselý}
\TULsupervisor{Prof. Ing. Zdeněk Molnár, CSc.}
\TULyear{2017}


%%%%%%%%%%%%%%%%%%%%%%%%%%%%
\begin{document}
\ThesisStart{male}


%%%%%%%%%%%%%%%%%%%%%%%%%%%%
\begin{abstractCZ}
Tato diplomová práce popisuje návrh a tvorbu fulltextového vyhledávání poskytovaného jako služba.
\end{abstractCZ}

\begin{klicovaslovaCZ}
Fulltext, Elasticsearch
\end{klicovaslovaCZ}

\vspace{2cm}

\begin{abstractEN}
This diploma thesis describes creation of fulltext search service.
\end{abstractEN}

\begin{klicovaslovaEN}
Fulltext, Elasticsearch
\end{klicovaslovaEN}


%%%%%%%%%%%%%%%%%%%%%%%%%%%%
\begin{acknowledgement}
Rád bych poděkoval Ing. Ivanovi Jelínkovi za rady a pomoc při řešení.
\end{acknowledgement}


%%%%%%%%%%%%%%%%%%%%%%%%%%%%
\tableofcontents
\clearpage


%%%%%%%%%%%%%%%%%%%%%%%%%%%%
\begin{abbrList}
\textbf{fulltextové vyhledávání} & 
Vyhledávání v indexovaných dokumentech, v české literatuře často označováno jako plnotextové vyhledávání\\
\textbf{e-shop} & 
Elektronický obchod\\
\textbf{e-commerce} & 
Elektronické obchodování, souhrnné označení obchodní činnosti prováděné na internetu\\
\end{abbrList}


%%%%%%%%%%%%%%%%%%%%%%%%%%%%
\chapter{Úvod}

V~této diplomové práci práci se zabývám návrhem a implementací nástroje, který umožní
provozovatelům e~shopů snadnou a rychlou implementaci fulltextového vyhledávání. Při
implementaci fulltextového vyhledávání je totiž třeba mít určitou znalost principů 
samotného vyhledávání, souvisejících oborů a technologií. Implementace kvalitního vyhledávání
tak může být pro provozovatele e-shopu náročná jak časově, tak finančně.

V rámci své diplomové práce navrhuji a vytvářím nástroj, který řeší tyto problémy a 
usnadňuje tak provozovatelům e-shopů spuštění vyhledávání. Konkrétní cíle diplomové práce
popisuji v následující kapitole. V další kapitole volím metody, kterými dané cíle naplním.

\section{Cíle práce}

Cílem práce je vytvořit nástroj, který bude nabízen jako služba a umožní rychlou implementaci
fulltextového vyhledávání do e-shopu. K~dosažení tohoto cíle je třeba naplnit cíle dílčí.

Prvním dílčím cíem je provedení analýzy problému a to jak z~pohledu byznysového, tak
z pohledu samotné problematiky vyhledávání. Co se týče pohledu byznysu, tak bude třeba
analyzovat potřeby zákazníků, zjistit, co je vlastně třeba vytvořit. Z~pohledu
problematiky vyhledávání je pak třeba provést analýzu této problematiky a utřídit
tak znalosti potřebné k poskytnutí kvalitních výsledků vyhledávání.

Dalším dílčím cílem je vytvořit návrh řešení problému na základě znalosti potřeb
zákazníků a dále na základě znalosti problematiky vyhledávání. Tento návrh bude v~dalším 
kroku implementovatelný.

Posledním dílčím cílem je implementace dle návrhu. Obnáší to výběr vhodných nástrojů
a konečně také ověření samotnoé implementace. Tu je třeba porovnat vůči požadavkům, 
aby se ověřilo, že byly splněny.

\section{Cílová skupina}

První cílovou skupinou je vývojář, řešící problém fulltextového vyhledávání produktů 
při implementaci elektronického obchodu. Takovému čtenáři by měla práce poskytnout dostatečné
teoretické znalosti potřebné pro implementaci vyhledávání. Užitečná také může být konkrétní
implementace, popisovaná v této práci.

Další možnou cílovou skupinou je provozovatel e-shopu, který přemýšlí, jakým způsobem
zlepšit fulltextové vyhledávání. V~této práci získá přehled o~složitosti samotné implementace,
poskytne mu komentovaný soupis možných řešení problémů a případně poskytne funkční
službu, kterou může okamžitě na~e-shop napojit.

Poslední cílovou skupinou je kdokoli, kdo se zajímá o fulltextové vyhledávání a českém jazyce.
V této práci nalezne soupis problémů souvisejicích s češtinou, které je třeba řešit. 
Dále zde nalezne konkrétní implementaci, kterou se může inspirovat při řešení obdobného problému.

\section{Použité metody}

V~této kapitole popisuji metody použité k~naplnění jednotlivých cílů. Pro zkoumání problému
oboru provedu analýzu trhu, nabízí se také možnost dotazování potenciálních uživatelů.
Pro porozumnění problematice vyhledávání provedu rešerši literatury. Vzhledem k množství
dostupné literatury bude třeba provést syntézu těchto informací. Pro vytvoření 
návrhu řešení bude použito ... Samotné vytvoření aplikace bude implementováno(?).

\section{Struktura práce}

V první části práce popisuji problematiku vyhledávání v~prostředí e-commerce. Definuji zde 
jednak kontext, v~kterém se pohybuji a popisuji problémy, které v~tomto prostředí řeším.
Snažím se identifikovat potencionálního zákazníka a definovat jeho požadavky na~vyhledávání.
Toto prostředí má svá specifika, která také popisuji. V~závěru této části porovnávám
existující nástroje umožňující implementaci vyhledávání a zjišťuji tak, proč má smysl
vytvářet další službu, jaká je její přidaná hodnota.

V druhé části se zabývám samotnou problematikou fulltextového vyhledávání. Popisuji zde celý
proces od analýzy vstupních dat, přes jejich indexaci, až po samotné vyhledávání. 
Tyto poznatky budou následně využity k~naplnění požadavků na vyhledávání, k~zajištění
kvalitních výsledků.

V další části práce navrhuji samotnou aplikaci, aby vyhovovala požadavkům a zároveň byla 
následně implementovatelná. Porovnávám zde dostupné nástroje, definuji příday užití aplikace
a vytvářím model nástroje.

V poslední části popisuji konkrétní implementaci v~jazyce Go s~pomocí databáze Elasticsearch.
Výstupem této části je otestovaný spustitelný program, který umožňuje provádět indexaci produktů
a jejich následné vyhledávání.


%%%%%%%%%%%%%%%%%%%%%%%%%%%%
%\chapter{Komentovaná rešerše informačních zdrojů}
%...


%%%%%%%%%%%%%%%%%%%%%%%%%%%%
\chapter{Analýza byznys požadavků na aplikaci}

\section{Vysvětlení problému, který je řešen}

\section{Specifika prostředí e-commerce}
- Specifika prostredi e-commerce, definice kontextu, ve kterem bude nastroj pouzivan

\section{Definice požadavků na vyhledávání}
- Definice pozadavku na samotne vyhledavani, akceptacni kriteria

\section{Technické požadavky na aplikaci}
- Technicke pozadavky (datove formaty, rychlost odezvy...)

\section{Popis oborů, kterých se práce dotýká}
- Ziskavani informaci, NLP, big data...

\section{Existující služby umožnující implementaci vyhledávání}
- Algolia, AWS CloudSearch...


%%%%%%%%%%%%%%%%%%%%%%%%%%%%
\chapter{Teorie fulltextového vyhledávání}

\section{Indexace dokumentů}

\section{Analýza dat, v kterých bude vyhledáváno}

\section{Zpracování přirozeného jazyka}
- NLP\\
- Sklonovani, casovani, tvaroslovi obecne\\
- Vaha slov, stop slova\\
- Preklepy\\
- Synonyma, zargon\\
- Lemmatizace\\
- Tezaurus\\
- ...

\chapter{Návrh řešení}

\section{Nástroje použitelné pro fulltextové vyhledáváni}
- Popis a porovnani vhodnych nastroju\\
\hspace*{5mm}- MySQL\\
\hspace*{5mm}- Elasticsearch\\
\hspace*{5mm}- Sphinx\\
\hspace*{5mm}- ...

\section{Definice případů užití}
- Uzivatel zada XML\\
- Uzivatel vyhleda
- ...

\section{Návrh architektury aplikace}
- Frontend vs backend (API)...


%%%%%%%%%%%%%%%%%%%%%%%%%%%%
\chapter{Implementace}
- Vyber technologii a nastroju (jazyk, framework...)

\section{Nastavení a nasazení Elasticsearch}
- Implementace indexace\\
- Implementace vyhledavani

\section{Backend}
- Popis trid, rozhrani, implementacni detaily

\section{Frontend}
- Frontend: UI, komunikace s backendem

\section{Deployment}
- Deployment, CI

\section{Testování, ověření}
- Overeni funkcnosti a prinosu\\
- Testovani, zda jsou splnena akceptacni kriteria\\
- Porovnani s stavajicim resenim vyhledavani na konkretnim priklade


%%%%%%%%%%%%%%%%%%%%%%%%%%%%
\chapter{Závěr}
V této diplomové práci jsem...

\section{Dosažení vytčených dílů}
...

\section{Diskuze možného budoucího rozšiřování}
...


%%%%%%%%%%%%%%%%%%%%%%%%%%%%
\begin{thebibliography}{Mm99}

\bibitem{strossa} STROSSA, Petr. \emph{Počítačové zpracování přirozeného jazyka}. \\
Vyd.~1. Praha: Oeconomica, 2011 316~s. ISBN 978-80-245-1777-3.

\bibitem{searching} AYSE Göker a DAVIES John. \emph{Information Retrieval: Searching in the 21st Century}. \\
Vyd.~1. Library of Congress Cataloging-in-Publication Data, 2009 295~s. ISBN: 978-0-470-02762-2.

\bibitem{mining} AGGARWAL Charu C., ZHAI ChengXiang. \emph{Mining Text Data}. \\
Springer New York Dordrecht Heidelberg London, 2000 522~s. ISBN 978-1-4614-3222-7.

\bibitem{es-guide} CLINTON Gormley, ZACHARY Tong. \emph{Elasticsearch: The Definitive Guide}. \\
O'Reilly Media, 2015 724~s. ISBN 978-1-4493-5854-9.

\bibitem{go-in-action} KENNEDY William. \emph{Go in Action}. \\
Manning Publications Co, 2016 264~s. ISBN 978-1-6172-9178-4.

\bibitem{elastic-reference} Elasticsearch. \emph{Elasticsearch Reference} [online]. \\
2017 [cit. 2017-03-03]. Dostupné z:\\
\url{https://www.elastic.co/guide/en/elasticsearch/reference/current/index.html}

\end{thebibliography}


%%%%%%%%%%%%%%%%%%%%%%%%%%%%
\chapter*{Rejstřík}
...


%%%%%%%%%%%%%%%%%%%%%%%%%%%%
\appendix

\chapter{Obsah přiloženého DVD}

\begin{itemize}
\item Soubor Diplomova\_prace\_2017\_Ludek\_Vesely.pdf
\begin{itemize}
	\item Text diplomové práce
\end{itemize}

\end{itemize}


%\chapter{Graf jako příloha}
%\begin{figure}[h]
%\center
%\includegraphics[width=0.86\textwidth]{todo.pdf}
%\end{figure}


%%%%%%%%%%%%%%%%%%%%%%%%%%%%
\end{document}
